\chapter*{DANH MỤC CÁC TỪ VIẾT TẮT VÀ THUẬT NGỮ}
\addcontentsline{toc}{chapter}{DANH MỤC CÁC TỪ VIẾT TẮT VÀ THUẬT NGỮ}

\begin{center}
\begin{tabular}{p{3cm} p{11cm}}
    \toprule
    \textbf{Ký hiệu/Viết tắt} & \textbf{Ý nghĩa đầy đủ / Thuật ngữ} \\
    \midrule
    
    AdamW     & Biến thể của thuật toán tối ưu Adam có tích hợp cơ chế phân rã trọng số (weight decay) tách biệt \\
    Aerosol sulfate & Các hạt sulfate ($SO_4^{2-}$) lơ lửng trong không khí, hình thành từ phản ứng của $SO_2$ với hơi nước trong điều kiện độ ẩm cao, là nguồn tạo $PM_{2.5}$ thứ cấp \\
    AI        & Artificial Intelligence (Trí tuệ nhân tạo) \\
    AQI       & Air Quality Index (Chỉ số chất lượng không khí) \\
    ARIMA     & Autoregressive Integrated Moving Average (Mô hình tự hồi quy tích hợp trung bình trượt) \\
    Batch size & Kích thước lô mẫu được nạp vào mô hình trong mỗi bước cập nhật tham số \\
    BiLSTM    & Bidirectional Long Short-Term Memory (LSTM hai chiều) \\
    Boosting  & Phương pháp học kết hợp tuần tự, mỗi mô hình mới sửa sai số của mô hình trước đó \\
    CART      & Classification and Regression Trees (Cây phân loại và hồi quy) \\
    CMAQ      & Community Multiscale Air Quality Model (Mô hình chất lượng không khí đa quy mô) \\
    CNN       & Convolutional Neural Network (Mạng thần kinh nhân chập) \\
    CO        & Carbon Monoxide (Khí carbon monoxit) \\
    COPD      & Chronic Obstructive Pulmonary Disease (Bệnh phổi tắc nghẽn mãn tính) \\
    CosineAnnealingLR & Bộ lịch trình giảm tốc độ học theo hàm cosin \\
    DL        & Deep Learning (Học sâu) \\
    Dropout   & Kỹ thuật chính quy hóa vô hiệu hóa ngẫu nhiên một tỷ lệ nơ-ron trong quá trình huấn luyện \\
    Early stopping & Cơ chế ngắt huấn luyện sớm khi hiệu suất trên tập kiểm định không cải thiện \\
    Ensemble  & Phương pháp kết hợp dự đoán từ nhiều mô hình thành phần nhằm nâng cao hiệu suất tổng thể \\
    Epoch     & Một vòng lặp hoàn chỉnh duyệt qua toàn bộ tập dữ liệu huấn luyện \\
    E-STGCN   & Enhanced Spatio-Temporal Graph Convolutional Network \\
    GCN       & Graph Convolutional Network (Mạng tích chập đồ thị) \\
    GELU      & Gaussian Error Linear Unit (Hàm kích hoạt tuyến tính lỗi Gaussian) \\
    Grid Search & Phương pháp tìm kiếm dạng lưới dò quét toàn bộ tổ hợp siêu tham số khả dĩ \\
    GRU       & Gated Recurrent Unit (Đơn vị hồi quy có cổng) \\
    Haversine & Công thức tính khoảng cách giữa hai điểm trên mặt cầu dựa trên tọa độ kinh vĩ độ \\
    Hessian   & Ma trận đạo hàm riêng bậc hai của hàm mục tiêu \\
    IoT       & Internet of Things (Internet vạn vật) \\
    IQR       & Interquartile Range (Khoảng tứ phân vị) \\
    iTransformer & Inverted Transformer --- kiến trúc Transformer đảo ngược chiều attention sang chiều biến \\
    KNN       & K-Nearest Neighbors (K láng giềng gần nhất) \\
    LightGBM  & Light Gradient Boosting Machine (Thuật toán tăng cường gradient nhẹ) \\
    LSTM      & Long Short-Term Memory (Bộ nhớ dài hạn ngắn hạn) \\
    MAE       & Mean Absolute Error (Sai số tuyệt đối trung bình) \\
    MAPE      & Mean Absolute Percentage Error (Sai số phần trăm tuyệt đối trung bình) \\
    ML        & Machine Learning (Học máy) \\
    MLP       & Multi-Layer Perceptron (Perceptron đa tầng) \\
    $NO_2$    & Nitrogen Dioxide (Khí nitơ dioxit) \\
    $O_3$     & Ozone (Khí ozon) \\
    OPC       & Optical Particle Counter (Bộ đếm hạt quang học) \\
    Overfitting & Hiện tượng quá khớp khi mô hình học thuộc nhiễu thay vì quy luật tổng quát \\
    Patience  & Số epoch liên tiếp không cải thiện trước khi kích hoạt early stopping \\
    $PM_{2.5}$  & Particulate Matter 2.5 (Bụi mịn có đường kính $\leq$ 2.5 $\mu m$) \\
    $PM_{10}$   & Particulate Matter 10 (Bụi có đường kính $\leq$ 10 $\mu m$) \\
    $R^2$     & Coefficient of Determination (Hệ số xác định) \\
    ReLU      & Rectified Linear Unit (Hàm kích hoạt tuyến tính chỉnh lưu) \\
    RevIN     & Reversible Instance Normalization (Chuẩn hóa mẫu khả nghịch) \\
    RMSE      & Root Mean Square Error (Sai số bình phương trung bình gốc) \\
    RNN       & Recurrent Neural Network (Mạng thần kinh hồi quy) \\
    ROS       & Reactive Oxygen Species (Các gốc oxy phản ứng) \\
    SARIMA    & Seasonal ARIMA (ARIMA có thành phần mùa vụ) \\
    Sigmoid   & Hàm kích hoạt dạng chữ S ánh xạ đầu ra về khoảng $[0, 1]$ \\
    SOA       & Secondary Organic Aerosol (Sol khí hữu cơ thứ cấp) \\
    $SO_2$    & Sulfur Dioxide (Khí lưu huỳnh dioxit) \\
    Streamlit & Nền tảng mã nguồn mở xây dựng giao diện Web cho ứng dụng khoa học dữ liệu \\
    ST-XLinear & Spatio-Temporal XLinear --- mô hình do đồ án đề xuất tích hợp đồ thị động vào XLinear \\
    SVM       & Support Vector Machine (Máy vector hỗ trợ) \\
    VOCs      & Volatile Organic Compounds (Các hợp chất hữu cơ dễ bay hơi) \\
    Weight decay & Hệ số phân rã trọng số trong thuật toán tối ưu, kiểm soát mức chính quy hóa L2 \\
    WHO       & World Health Organization (Tổ chức Y tế Thế giới) \\
    WRF       & Weather Research and Forecasting Model (Mô hình nghiên cứu và dự báo thời tiết) \\
    XGBoost   & Extreme Gradient Boosting \\
    XLinear   & Mô hình dự báo chuỗi thời gian dựa trên MLP với cơ chế Gating Block cho biến ngoại sinh \\
    
    \bottomrule
\end{tabular}
\end{center}
